\subsection{A threefold with a nonzero integral restriction kernel}

\begin{lemma}[The positive-section Tango construction, with its signs fixed]
\label{lem:OC-positive-Tango-iteration}
\status{Proved; independently audited by the root and reduction-descent agents.}
Let $B$ be a smooth projective connected variety over an algebraically closed field of characteristic two.  Suppose that an effective ample divisor $D$ and a rational function $f$ satisfy the following stronger form of the Tango condition: $df$, divided locally by an equation of $2D$, is a regular nowhere-zero differential.  Let $k>2$ be odd, suppose $D=kN$ for an effective divisor $N$, and put $T=\OO_B(D)$.
There is a smooth projective connected variety $\widetilde B$ of dimension $\dim B+1$, a morphism $g:\widetilde B\to B$, and a smooth section $E$ of $g$, with the following properties:
\begin{enumerate}
\item If $\pi:P=\mathbf P_B(V)\to B$ parametrizes quotients of the extension
\[
 0\longrightarrow\OO_B\longrightarrow V\longrightarrow T\longrightarrow0
\]
defined by $df$, and $E_0$ is its quotient section, there is a finite flat normalized cyclic cover $\psi:\widetilde B\to P$ of degree $k$, with $g=\pi\psi$ and $\psi^*E_0=kE$.
\item For every integer $a$, its normalized eigensheaves are
\begin{equation}\label{eq:OC-iteration-eigensheaves}
 \psi_*\OO_{\widetilde B}(aE)
 =\bigoplus_{i=0}^{k-1}
   \OO_P\!\left(\left\lfloor\frac{a-2i}{k}\right\rfloor\right)
                       \otimes\pi^*\OO_B(2iN).
\end{equation}
\item The divisor $\widetilde D=(k-1)E+g^*N$ is ample and effective, and a rational function $\widetilde f$ on $\widetilde B$ satisfies the same nowhere-zero Tango condition with $\widetilde D$.
\item For every prime $\ell\ne2$, pullback identifies
$H^1_{\et}(B,\mathbf Q_\ell)$ with $H^1_{\et}(\widetilde B,\mathbf Q_\ell)$.
\end{enumerate}
\end{lemma}
\begin{proof}
The exact sequence for the sheaf of exact differentials gives a nonzero class
\[
 \ker\{H^1(B,T^{-1})\xrightarrow{F}H^1(B,T^{-2})\};
\]
this is \cite[Theorem~1, Lemma~1.4, and their proofs, pp.~1, 4--5]{Mukai2011Kodaira}.
It defines the stated extension, split after Frobenius pullback.  We derive the covering signs in the quotient-projective-bundle convention instead of importing the sign in the printed formula for the cyclic line bundle.

Choose local equations $b_i$ of $N$ and $a_i=b_i^k$ of $D$.  By the Cartier exact sequence, after refining the cover there are regular functions $f_i$ and rational functions $h_i$ with
\[
       f/a_i^2=f_i+h_i^2,
       \qquad df_i=a_i^{-2}df,
\]
where $df_i$ is nowhere zero.  On overlaps set $t_{ij}=a_i/a_j$.  Then
\[
 f_j=t_{ij}^2f_i+c_{ij}^2,
 \qquad c_{ij}=t_{ij}h_i-h_j.
\]
The $c_{ij}$ are regular: their squares are regular and the overlap rings are normal.  The affine coordinates in $P\setminus E_0$ glue by
$S_j=t_{ij}S_i+c_{ij}$.  This is the Frobenius-split extension construction; its infinity section has normal bundle $T$ and
\[
 \OO_P(E_0)=\OO_P(1).
\]
The equations $S_i^2=f_i$ define a multisection $G$ disjoint from $E_0$.  It is smooth because $df_i$ is nowhere zero, and its divisor bundle is
\[
             \OO_P(G)=\OO_P(2)\otimes\pi^*T^{-2}.
\]
These local calculations also verify directly the geometry used in
\cite[Proposition~1.7 and its proof, pp.~5--6]{Mukai2011Kodaira}.

Set $m=k-2$ and
\[
 M=\OO_P(-1)\otimes\pi^*\OO_B(2N),
 \qquad M^{-k}=\OO_P(G+mE_0).
\]
Normalize the cyclic cover determined by this branch divisor.  At $G$ the local equation has a simple branch; at $E_0$, after an \'etale extension extracting a root of a unit, it is $w^k=s^{k-2}$.  Since $\gcd(k,k-2)=1$, the normalization there is $s=u^k$, $w=u^{k-2}$.  The two branch divisors are disjoint and smooth, so these equations prove that the normalization is smooth.  The valuation one at the integral divisor $G$ makes the generic cyclic polynomial Eisenstein; thus the cover is integral.

The index-$i$ normalization allows the pole
$\lfloor(mi+a)/k\rfloor$ along $E_0$.  Its eigensheaf is
\[
 M^i\bigl(\lfloor(mi+a)/k\rfloor E_0\bigr)
 =\OO_P\!\left(\left\lfloor\frac{a-2i}{k}\right\rfloor\right)
                      \otimes\pi^*\OO_B(2iN).
\]
This proves \eqref{eq:OC-iteration-eigensheaves}; in particular the normalized algebra is locally free, proving flatness.  It is the normalized construction of \cite[Definition~2.1, p.~7]{Mukai2011Kodaira}, expressed in our explicitly checked positive-section convention.

On the affine chart write
\[
                         T_i^k=S_i^2-f_i.
\]
The gluing $T_j=(b_i/b_j)^2T_i$ makes
$\widetilde f=b_i^2T_i$ a global rational function.  Its differential is $b_i^2dT_i$.  Off $G$ and $E$, the cover is \'etale and
\[
                   kT_i^{k-1}dT_i=-df_i.
\]
At $G$, $f_i$ is part of an \'etale coordinate system on $B$, so the equation solves for that coordinate and $T_i$ is a smooth local parameter; hence $dT_i$ is nowhere zero there.  Along $E$, $T_i$ has pole order two.  The same differential identity then shows that $dT_i$ has exact zero order $2(k-1)$: the pullback of $df_i$ is nowhere zero along $E$, since $E\simeq B$ and the map is smooth in its normal direction there.  Thus
\[
       (d\widetilde f)=2\bigl((k-1)E+g^*N\bigr)
\]
in the precise sense of a nowhere-zero normalized differential.  This supplies directly the iteration asserted in
\cite[Proposition~2.3 and its proof, p.~8]{Mukai2011Kodaira}.

Every geometric fibre of $g$ has affine equation $T^k=(S-c)^2$ and one smooth point at infinity.  It is integral, with normalization $\mathbf P^1$, and the normalization is bijective and radicial.  The divisor $E$ has degree one on every fibre.  Moreover
\[
       (\OO(E)|_E)^k\simeq\OO_B(D)=\OO_B(kN),
\]
so its restriction to $E$ is numerically ample.  The effective divisor $E$ therefore has nonnegative degree on every curve, including curves contained in $E$; it is nef.  For an irreducible subvariety $Z\subset\widetilde B$ of dimension $d$, its image under $g$ has dimension $d$ or $d-1$.  In the first case $(g^*N)^d\cdot Z>0$; in the second case $E(g^*N)^{d-1}\cdot Z>0$, by the positive fibre degree.  All other mixed terms are nonnegative.  Nakai--Moishezon now proves that $(k-1)E+g^*N$ is ample.

Finally the finite normalization of a geometric fibre has one-point geometric fibres.  Proper base change therefore identifies its constant-sheaf cohomology with that of $\mathbf P^1$, so its $H^1_{\et}$ is zero.  Proper base change again gives $g_*\mathbf Q_\ell=\mathbf Q_\ell$ and $R^1g_*\mathbf Q_\ell=0$, and the Leray edge gives the asserted $H^1$ isomorphism.  One may apply this first to $\mathbf Z/\ell^n$ and pass to the usual inverse limit.  The exact source, inspected in a retypesetting of the original text with retained original pagination, is \cite[Expos\'e~XII, Theorem~5.1 and Corollary~5.2, original pp.~88--89]{SGA4III1973}.
\end{proof}

\begin{proposition}[An explicit generated Tango threefold]
\label{prop:OC-Tango-threefold}
\status{Proved; independently audited by the root and reduction-descent agents.}
In characteristic two there is an explicitly specified smooth projective threefold $R$ and an ample globally generated line bundle $L$ such that
\[
 H^1(R,L^{-2})\ne0,\qquad
 L^{445}\ \hbox{is very ample},\qquad
 \dim H^1_{\et}(R,\mathbf Q_3)=842.
\]
The construction starts with $C:y^2+y=x^{843}$ and applies Lemma~\ref{lem:OC-positive-Tango-iteration} first with $k=21$ and then with $k=5$.
\end{proposition}
\begin{proof}
Let $C$ be the smooth projective model of the displayed Artin--Schreier curve.  The same direct local computation as in Lemma~\ref{lem:OC-generated-Tango-root}, now with $h=422$, gives
\[
 g(C)=421,\qquad (dx)=840P_\infty,\qquad
 N_0=\OO_C(2P_\infty),\qquad D_1=420P_\infty.
\]
For completeness, $v(x)=-2$, $v(y)=-843$, $w=y/x^{422}$ is a parameter at infinity, and $v=1/x$ satisfies $v+w v^{422}=w^2$.  The formula for $dx$ gives order $2(422-2)=840$, with no other zero or pole.  Thus $D_1=210N_0$, and $N_0$ is generated by $1,x$.

For the first step take $k=21$ and $N=10N_0$, obtaining $f:S\to C$ and the ramification section $E_1$.  Its new Tango divisor is
\[
             D_2=20E_1+10f^*N_0=10L_2,
 \qquad L_2=\OO_S(2E_1)\otimes f^*N_0,
\]
where equality with a line bundle denotes the associated divisor class and the effective representative just displayed is fixed.  Since $D_2$ is ample, $L_2$ is ample.  Formula~\eqref{eq:OC-iteration-eigensheaves} at $a=2$ gives
\[
                         f_*L_2=N_0\oplus N_0^{21}.
\]
Indeed only indices $i=0,1$ have relative degree zero; the remaining degrees are negative.  The index-zero sections generate away from $E_1$, and the index-one local generator is $w/s=u^{-2}$ along $E_1$.  The two coefficient bundles are generated, proving that $L_2$ is generated everywhere.

We give the positivity bounds needed for the next step.  On the first projective bundle the extension is $0\to\OO_C\to V_1\to N_0^{210}\to0$.  The bundle $\operatorname{Sym}^aV_1\otimes N_0^b$ has a filtration with line quotients $N_0^{b+210j}$, $0\leq j\leq a$.  If $a\geq1$ and $b\geq441$, all quotient degrees are at least $882>2g(C)+1=843$.  Serre duality gives vanishing of $H^1$ after subtracting any length-two subscheme of $C$.  Global sections therefore supply all values and first jets on the base, as well as the complete fibre Veronese system.  Consequently
\[
                   \OO_{\mathbf P(V_1)}(a)\otimes N_0^b
                   \quad\hbox{is very ample}.
\]
Now $L_2^{441}=\psi_1^*F_1$ for
$F_1=\OO_{\mathbf P(V_1)}(42)\otimes N_0^{441}$.
The normalized cover algebra has summands
\[
 K_i=\OO_{\mathbf P(V_1)}\!\left(-\left\lceil\frac{2i}{21}\right\rceil\right)
                      \otimes N_0^{20i},\qquad 0\leq i\leq20.
\]
All $F_1\otimes K_i$ are globally generated, since their relative degrees are at least $40$ and their base exponents at least $441$.  A finite-cover coordinate argument proves that $\psi_1^*F_1$ is very ample: sections in the index-zero summand recover the embedded base, and sections generating every $F_1\otimes K_i$ recover the entire finite algebra on each affine chart.  Hence $L_2^{441}$ is very ample.

For every integer $t\geq441$, the eigensheaves of $L_2^t$ push to
\[
 \operatorname{Sym}^{a_i(t)}V_1\otimes N_0^{t+20i},
 \qquad a_i(t)=\left\lfloor\frac{2(t-i)}{21}\right\rfloor\geq40.
\]
Their filtered line quotients all have degree at least $2t\geq882>840=2g(C)-2$.  Their $H^1$ groups vanish, and the nonnegative relative degrees give no higher direct image.  Thus
\begin{equation}\label{eq:OC-L2-jet-bound}
 H^1(S,L_2^t)=0,\qquad
 H^1(S,L_2^t\otimes I_Z)=0
 \quad(t\geq441,\ \operatorname{length}Z\leq2).
\end{equation}
For the second assertion, $L_2^t$ is very ample because it is $L_2^{441}$ tensored with a generated power; combine separation of length-two schemes with the first vanishing and its evaluation sequence.

Apply the Tango step again, to $D_2=10L_2$, now with $k=5$ and $N=2L_2$.  It gives $g:R\to S$, section $E_2$, and
\[
              D_3=4E_2+2g^*L_2=2L,
 \qquad L=\OO_R(2E_2)\otimes g^*L_2.
\]
As before $L$ is ample, $g_*L=L_2\oplus L_2^5$, and the index-zero and index-one local generators prove global generation.  The nowhere-zero exact differential with divisor $2D_3$ gives a nonzero class in
$H^1(R,\OO_R(-D_3))=H^1(R,L^{-2})$ by the exact-differential sequence cited in the preceding lemma.

On the second projective bundle $P_2=\mathbf P_S(V_2)$ the extension is $0\to\OO_S\to V_2\to L_2^{10}\to0$.  If $a\geq1$ and $b\geq441$, the filtration of $\operatorname{Sym}^aV_2\otimes L_2^b$ has quotients $L_2^{b+10j}$; \eqref{eq:OC-L2-jet-bound} makes its restriction to any length-two subscheme of $S$ surjective on global sections.  This supplies horizontal jets, while $\OO_{P_2}(a)$ separates all fibre jets.  Thus $\OO_{P_2}(a)\otimes L_2^b$ is very ample.  More explicitly, a length-two subscheme of $P_2$ either projects isomorphically to a length-two subscheme of $S$, or is supported in a fibre; the two preceding surjections handle these cases, including a tangent direction with nonzero base component.

Because $445=5\cdot89$,
\[
 L^{445}=\psi_2^*F_2,
 \qquad F_2=\OO_{P_2}(178)\otimes L_2^{445}.
\]
The normalized algebra summands are
\[
 K_j'=\OO_{P_2}\!\left(-\left\lceil\frac{2j}{5}\right\rceil\right)
                              \otimes L_2^{4j},\qquad0\leq j\leq4.
\]
Every $F_2\otimes K_j'$ has relative degree at least $176$ and base exponent at least $445$.  They are globally generated and $F_2$ is very ample.  The same finite-algebra coordinate argument proves that $L^{445}$ is very ample.

The final assertion follows twice from the last part of Lemma~\ref{lem:OC-positive-Tango-iteration}:
\[
       H^1_{\et}(R,\mathbf Q_3)\simeq H^1_{\et}(S,\mathbf Q_3)
                      \simeq H^1_{\et}(C,\mathbf Q_3),
\]
whose dimension is $2g(C)=842$.
\end{proof}

\begin{lemma}[A multiroot cover amplifies coherent kernels and preserves first Betti number]
\label{lem:OC-threefold-multiroot}
\status{Proved; independently audited by the root and reduction-descent agents.}
For $(R,L)$ in Proposition~\ref{prop:OC-Tango-threefold}, choose $843$ sufficiently general sections $b_i$ of $L^{445}$, of which the first eight define a closed immersion into $\mathbf P^7$.  The simultaneous root cover
\[
 \rho:X=\operatorname{Spec}_R
  \left(\bigotimes_{i=1}^{843}\bigoplus_{a_i=0}^{444}L^{-a_i}\right)
       \longrightarrow R,
 \qquad y_i^{445}=b_i,
\]
is smooth projective connected of degree $445^{843}$, the line bundle $A=\rho^*L$ is very ample, and
\[
 h^1(X,A^{-1})\geq843,
 \qquad \dim H^1_{\et}(X,\mathbf Q_3)=842.
\]
\end{lemma}
\begin{proof}
The first eight sections exist because $L^{445}$ is very ample on a threefold.  Indeed the closure of the union of its secant and tangent lines has dimension at most seven.  A generic linear-projection centre of codimension eight misses this closed subset; projection to $\mathbf P^7$ therefore separates all length-two subschemes and is a closed immersion.  Here the index-zero summand pushes to $\operatorname{Sym}^{178}V_2\otimes L_2^{445}$, whose $179$ filtered line quotients have vanishing $H^1$ and at least one global section each.  Thus $h^0(R,L^{445})\geq179>8$, and the projection has the indicated dimension.  The accompanying exact Riemann--Roch certificate gives the stronger count $79{,}232{,}578{,}014$.  The independence and immersion conditions form a nonempty open subset of the space of eight-tuples.

Simultaneously choose all $843$ sections so that their zero divisors $D_i$ are smooth irreducible surfaces meeting with simple normal crossings, with at most three divisors through any point.  These are also nonempty open conditions: apply Bertini to the successive smooth intersections and avoid the finitely many existing triple-intersection points.  Smooth ample divisors on the connected threefold are connected.  The parameter space of tuples is irreducible, so these conditions and the first-eight embedding condition hold together.

The cover algebra is locally free of rank $445^{843}$.  At a point on $a\leq3$ branch divisors their equations are part of a regular parameter system; adjoining their roots replaces these parameters by new parameters.  The other roots are \'etale extensions, so $X$ is smooth.  At the valuation of $D_i$, all other branch equations are units.  The earlier extensions in the tower are consequently unramified there, while the next equation $T^{445}-b_i$ is Eisenstein.  Its degree is $445$ at each stage.  Thus the generic algebra is a field, and flatness shows that the smooth cover is integral and connected.

Take any generating system $s_\alpha$ for $L$.  The sections $\rho^*s_\alpha$ and all $y_i$ define the required projective morphism.  On the chart $s_\alpha\ne0$, the functions $b_i/s_\alpha^{445}$ for the first eight indices generate the affine coordinate ring of $R_{s_\alpha}$: enlarge their already very ample linear system by the section $s_\alpha^{445}$ and take its corresponding affine chart.  Enlarging a very ample system preserves its closed immersion.  The ratios $y_i/s_\alpha$ then generate the entire cover algebra and their $445$th powers recover $b_i/s_\alpha^{445}$.  Hence the displayed projective morphism is a closed immersion on a covering of $X$.  Properness proves the global assertion.  Four general sections suffice to generate $L$ on a threefold, if one wants a definite list of $847$ embedding sections.

The eigensheaf decomposition and projection formula give
\[
 H^1(X,A^{-1})=
 \bigoplus_{0\leq a_i<445}
 H^1\!\left(R,L^{-(1+\sum a_i)}\right).
\]
There are $843$ indices with $\sum a_i=1$, each contributing the nonzero group $H^1(R,L^{-2})$.  This proves the lower bound.

For the Betti number, extend $\mathbf Q_3$ to a finite field $K$ containing the $445$th roots of unity.  Decompose $\rho_*K$ into the characters of $(\mu_{445})^{843}$.  The trivial character gives $K_R$.  For a nontrivial character $\chi$, let $J$ be the nonempty set of branch indices on whose inertia it is nontrivial, and let
\[
 U=R\setminus\bigcup_{i\in J}D_i,
 \qquad j:U\hookrightarrow R.
\]
Its summand is $j_!\mathcal L_\chi$, where $\mathcal L_\chi$ is a lisse rank-one sheaf on $U$.  This can be checked on each root-cover factor: a trivial character extends constantly across its branch, whereas a nontrivial character has zero stalk on its branch; their tensor product has exactly the asserted extension by zero.
The open set $U$ is affine, since the removed effective divisor is ample, and it is smooth of dimension three.  Artin affine vanishing and Poincar\'e duality yield
\[
 H^1(R,j_!\mathcal L_\chi)=H_c^1(U,\mathcal L_\chi)
   \simeq H^5(U,\mathcal L_\chi^\vee(3))^\vee=0.
\]
The exact primary statements are \cite[Expos\'e~XIV, Theorem~3.1 and Corollary~3.2, original p.~160; Expos\'e~XVIII, Theorem~3.2.5 and formula~(3.2.6.2), original pp.~585--587]{SGA4III1973}.  They apply first with finite coefficient lattices; $445$ is a unit at $3$, so the character decomposition and passage to $K$ preserve the assertion.  Finite proper base change gives no higher direct images for $\rho$.  All nontrivial characters therefore contribute zero to $H^1$, and
\[
                      H^1_{\et}(X,K)=H^1_{\et}(R,K).
\]
This proves the stated Betti number.
\end{proof}

\begin{theorem}[Integral crystalline hyperplane restriction need not be injective]
\label{thm:OC-integral-injectivity-counterexample}
\status{Proved; independently audited by the root and reduction-descent agents.}
In characteristic two, let $(X,A)$ be the explicit smooth threefold of Lemma~\ref{lem:OC-threefold-multiroot}.  For every smooth hyperplane section $Y\in|A|$,
\[
 \ker\bigl(H^2_{\cris}(X/W)\longrightarrow H^2_{\cris}(Y/W)\bigr)[2]
                         \ne0.
\]
Thus the universal injectivity assertion A.1(c) fails already in dimension three and degree one.
\end{theorem}
\begin{proof}
The smooth ample section $Y$ is connected.  The sequence
$0\to A^{-1}\to\OO_X\to\OO_Y\to0$ identifies its coherent restriction kernel with $H^1(X,A^{-1})$, which has dimension at least $843$.  The functorial conjugate edge injection
\[
           H^1(X^{(1)},\OO_{X^{(1)}})\hookrightarrow H^1_{\dR}(X/k)
\]
therefore puts a subspace of dimension at least $843$ inside
\[
       K=\ker\{H^1_{\dR}(X/k)\longrightarrow H^1_{\dR}(Y/k)\}.
\]
The Frobenius twist preserves dimension because the ground field is perfect.  The edge injection is the one explicitly displayed in
\cite[II, Section~6.14, p.~649]{Illusie1979DRW}; it also follows directly from Cartier and the degree-one conjugate spectral sequence.

The integral universal-coefficient sequence, natural for restriction, is
\[
 0\longrightarrow H^1_{\cris}(X/W)/2
 \longrightarrow H^1_{\dR}(X/k)
 \xrightarrow{\beta_X}H^2_{\cris}(X/W)[2]\longrightarrow0.
\]
The first crystalline group is free and has rank
$\dim H^1_{\et}(X,\mathbf Q_3)=842$: the exact source is
\cite[II, Proposition~3.11 and Remark~3.11.2, pp.~618--619]{Illusie1979DRW}, identifying it with the Dieudonn\'e realization of the reduced Picard variety, whose height is its first Betti number.  Consequently $K$ cannot be contained in the $842$-dimensional kernel of $\beta_X$.  Choose $v\in K$ with $\beta_X(v)\ne0$.  Naturality gives
\[
       \operatorname{res}(\beta_X(v))
         =\beta_Y(\operatorname{res}(v))=0.
\]
This is the required nonzero element of the integral restriction kernel, killed by two.
\end{proof}

\begin{remark}[Explicit choices and the exact degree refinement]
\label{rem:OC-injectivity-degrees}
\status{Proved under stated hypotheses}
The exact arithmetic certificate \texttt{A\_injectivity\_checks.py} verifies $52{,}026$ normalization-floor identities, both jet bounds, the Riemann--Roch counts, and the strict dimension inequality; its recorded run returned PASS in $0.013554183002270292$ seconds.  These finite checks supplement the geometric proofs and do not replace them.

All generic choices can be specified over the algebraic closure of a rational function field over $\overline{\mathbf F}_2$: take independent generic coefficient vectors in a chosen basis of $H^0(R,L^{445})$, then independent generic coefficients for the final hyperplane.  Replacing the initial basis merely makes an invertible change of those generic coordinates.  The fixed integers are the curve exponent $843$, the cyclic degrees $21,5$, the very ample power $445$, and the number $843$ of root equations.  The proof establishes degree-one failure for every smooth hyperplane in the resulting embedding.  The exact degree range is determined in Theorem~\ref{thm:OC-exact-integral-injectivity-degrees}: injectivity fails in degrees one and two and holds for every degree at least three.
\end{remark}
